\begingroup\raggedright\small
\begin{thebibliography}{CHOPRS22}
\bibitem[Aar16]{Aar16} S. Aaronson, \emph{$\mathsf P\overset{?}{=}\mathsf{NP}$}, in: Open Problems in Mathematics, J. F. Nash, Jr. and M. Th. Rassias (eds.), Springer, Cham, 2016, pp.~1--122. \href{https://doi.org/10.1007/978-3-319-32162-2_1}{doi:10.1007/978-3-319-32162-2\_1}
\bibitem[AB09]{AB09} S. Arora and B. Barak, \emph{Computational Complexity}, Cambridge University Press, Cambridge, 2009. \href{https://doi.org/10.1017/CBO9780511804090}{doi:10.1017/CBO9780511804090}
\bibitem[AB87]{AB87} N. Alon and R. B. Boppana, \emph{The monotone circuit complexity of boolean functions}, Combinatorica \textbf{7} (1987), no.~1, 1--22. \href{https://doi.org/10.1007/BF02579196}{doi:10.1007/BF02579196}
\bibitem[Ajt83]{Ajt83} M. Ajtai, \emph{$\Sigma^1_1$-formulae on finite structures}, Ann. Pure Appl. Logic \textbf{24} (1983), no.~1, 1--48. \href{https://doi.org/10.1016/0168-0072(83)90038-6}{doi:10.1016/0168-0072(83)90038-6}
\bibitem[AK10]{AK10} E. Allender and M. Kouck{\'y}, \emph{Amplifying lower bounds by means of self-reducibility}, J. ACM \textbf{57} (2010), no.~3, 1--36. \href{https://doi.org/10.1145/1706591.1706594}{doi:10.1145/1706591.1706594}
\bibitem[AKS04]{AKS04} M. Agrawal, N. Kayal and N. Saxena, \emph{PRIMES is in P}, Ann. of Math. (2) \textbf{160} (2004), no.~2, 781--793. \href{https://doi.org/10.4007/annals.2004.160.781}{doi:10.4007/annals.2004.160.781}
\bibitem[ALMSS98]{ALMSS98} S. Arora, C. Lund, R. Motwani, M. Sudan and M. Szegedy, \emph{Proof verification and the hardness of approximation problems}, J. ACM \textbf{45} (1998), no.~3, 501--555. \href{https://doi.org/10.1145/278298.278306}{doi:10.1145/278298.278306}
\bibitem[AS98]{AS98} S. Arora and S. Safra, \emph{Probabilistic checking of proofs}, J. ACM \textbf{45} (1998), no.~1, 70--122. \href{https://doi.org/10.1145/273865.273901}{doi:10.1145/273865.273901}
\bibitem[AW09]{AW09} S. Aaronson and A. Wigderson, \emph{Algebrization}, ACM Trans. Comput. Theory \textbf{1} (2009), no.~1, 1--54. \href{https://doi.org/10.1145/1490270.1490272}{doi:10.1145/1490270.1490272}
\bibitem[BGS75]{BGS75} T. Baker, J. Gill and R. Solovay, \emph{Relativizations of the $\mathcal P=?\mathcal{NP}$ question}, SIAM J. Comput. \textbf{4} (1975), no.~4, 431--442. \href{https://doi.org/10.1137/0204037}{doi:10.1137/0204037}
\bibitem[Bha26bsd]{Bha26bsd} D. Bhattacharjee, \emph{The $p$-primary part of the Tate--Shafarevich group in ranks three to six: paper, data and certificates}, Zenodo, release v1.0.2 (2026). \href{https://doi.org/10.5281/zenodo.23933270}{doi:10.5281/zenodo.23933270}
\bibitem[Bha26eth]{Bha26eth} D. Bhattacharjee, \emph{Source and verification code for the paper Subexponential Weighted Counting and Exact Optimisation}, Zenodo, release v1.4.0 (2026). \href{https://doi.org/10.5281/zenodo.23718560}{doi:10.5281/zenodo.23718560}
\bibitem[Bha26weil]{Bha26weil} D. Bhattacharjee, \emph{LaTeX source and verification code for the paper Algebraic Loci of Weil Classes on Abelian Varieties}, Zenodo, release v5.1.0 (2026). \href{https://doi.org/10.5281/zenodo.23227940}{doi:10.5281/zenodo.23227940}
\bibitem[BhaHex26]{BhaHex26} D. Bhattacharjee, \emph{Subexponential weighted counting and exact optimisation}, HexagonMath preprint 2610.00594v1 (2026), \url{https://hexagonmath.org/2610.00594v1}.
\bibitem[BIP19]{BIP19} P. B{\"u}rgisser, C. Ikenmeyer and G. Panova, \emph{No occurrence obstructions in geometric complexity theory}, J. Amer. Math. Soc. \textbf{32} (2019), no.~1, 163--193. \href{https://doi.org/10.1090/jams/908}{doi:10.1090/jams/908}
\bibitem[Blu84]{Blu84} N. Blum, \emph{A Boolean function requiring $3n$ network size}, Theoret. Comput. Sci. \textbf{28} (1984), no.~3, 337--345. \href{https://doi.org/10.1016/0304-3975(83)90029-4}{doi:10.1016/0304-3975(83)90029-4}
\bibitem[BSW01]{BSW01} E. Ben-Sasson and A. Wigderson, \emph{Short proofs are narrow---resolution made simple}, J. ACM \textbf{48} (2001), no.~2, 149--169. \href{https://doi.org/10.1145/375827.375835}{doi:10.1145/375827.375835}
\bibitem[Cai07]{Cai07} J.-Y. Cai, \emph{$\mathrm{S}_2^p\subseteq\mathrm{ZPP}^{\mathrm{NP}}$}, J. Comput. System Sci. \textbf{73} (2007), no.~1, 25--35. \href{https://doi.org/10.1016/j.jcss.2003.07.015}{doi:10.1016/j.jcss.2003.07.015}
\bibitem[Che11]{Che11} D. Yu. Cherukhin, \emph{Lower bounds for complexity of Boolean circuits of finite depth with arbitrary elements}, Discrete Math. Appl. \textbf{21} (2011), no.~4. \href{https://doi.org/10.1515/dma.2011.031}{doi:10.1515/dma.2011.031}
\bibitem[CHOPRS22]{CHOPRS22} L. Chen, S. Hirahara, I. C. Oliveira, J. Pich, N. Rajgopal and R. Santhanam, \emph{Beyond Natural Proofs: Hardness Magnification and Locality}, J. ACM \textbf{69} (2022), no.~4, 1--49. \href{https://doi.org/10.1145/3538391}{doi:10.1145/3538391}
\bibitem[CHR24]{CHR24} L. Chen, S. Hirahara and H. Ren, \emph{Symmetric Exponential Time Requires Near-Maximum Circuit Size}, in: Proceedings of the 56th Annual ACM Symposium on Theory of Computing (STOC 2024), 2024, pp.~1990--1999. \href{https://doi.org/10.1145/3618260.3649624}{doi:10.1145/3618260.3649624}
\bibitem[CIP06]{CIP06} C. Calabro, R. Impagliazzo and R. Paturi, \emph{A duality between clause width and clause density for SAT}, in: 21st Annual IEEE Conference on Computational Complexity (CCC 2006), 2006, pp.~252--260. \href{https://doi.org/10.1109/CCC.2006.6}{doi:10.1109/CCC.2006.6}
\bibitem[CIP09]{CIP09} C. Calabro, R. Impagliazzo and R. Paturi, \emph{The Complexity of Satisfiability of Small Depth Circuits}, in: Parameterized and Exact Computation (IWPEC 2009), Lecture Notes in Comput. Sci. \textbf{5917}, Springer, Berlin, 2009, pp.~75--85. \href{https://doi.org/10.1007/978-3-642-11269-0_6}{doi:10.1007/978-3-642-11269-0\_6}
\bibitem[CKKSZ15]{CKKSZ15} R. Chen, V. Kabanets, A. Kolokolova, R. Shaltiel and D. Zuckerman, \emph{Mining Circuit Lower Bound Proofs for Meta-Algorithms}, Comput. Complexity \textbf{24} (2015), no.~2, 333--392. \href{https://doi.org/10.1007/s00037-015-0100-0}{doi:10.1007/s00037-015-0100-0}
\bibitem[Clay]{Clay} Clay Mathematics Institute, \emph{P vs NP}, Millennium Prize Problems, official problem description by S. Cook, \url{https://www.claymath.org/millennium/p-vs-np/}.
\bibitem[Coo71]{Coo71} S. A. Cook, \emph{The complexity of theorem-proving procedures}, in: Proceedings of the Third Annual ACM Symposium on Theory of Computing (STOC 1971), 1971, pp.~151--158. \href{https://doi.org/10.1145/800157.805047}{doi:10.1145/800157.805047}
\bibitem[Coo73]{Coo73} S. A. Cook, \emph{A hierarchy for nondeterministic time complexity}, J. Comput. System Sci. \textbf{7} (1973), no.~4, 343--353. \href{https://doi.org/10.1016/S0022-0000(73)80028-5}{doi:10.1016/S0022-0000(73)80028-5}
\bibitem[CR79]{CR79} S. A. Cook and R. A. Reckhow, \emph{The relative efficiency of propositional proof systems}, J. Symbolic Logic \textbf{44} (1979), no.~1, 36--50. \href{https://doi.org/10.2307/2273702}{doi:10.2307/2273702}
\bibitem[CSS18]{CSS18} R. Chen, R. Santhanam and S. Srinivasan, \emph{Average-case lower bounds and satisfiability algorithms for small threshold circuits}, Theory Comput. \textbf{14} (2018), no.~1, 1--55. \href{https://doi.org/10.4086/toc.2018.v014a009}{doi:10.4086/toc.2018.v014a009}
\bibitem[Din07]{Din07} I. Dinur, \emph{The PCP theorem by gap amplification}, J. ACM \textbf{54} (2007), no.~3, Art.~12. \href{https://doi.org/10.1145/1236457.1236459}{doi:10.1145/1236457.1236459}
\bibitem[ECCC]{ECCC} Electronic Colloquium on Computational Complexity, \url{https://eccc.weizmann.ac.il/}.
\bibitem[Edm65]{Edm65} J. Edmonds, \emph{Paths, Trees, and Flowers}, Canad. J. Math. \textbf{17} (1965), 449--467. \href{https://doi.org/10.4153/CJM-1965-045-4}{doi:10.4153/CJM-1965-045-4}
\bibitem[FGHK16]{FGHK16} M. G. Find, A. Golovnev, E. A. Hirsch and A. S. Kulikov, \emph{A better-than-$3n$ lower bound for the circuit complexity of an explicit function}, in: 57th Annual IEEE Symposium on Foundations of Computer Science (FOCS 2016), 2016, pp.~89--98. \href{https://doi.org/10.1109/FOCS.2016.19}{doi:10.1109/FOCS.2016.19}
\bibitem[FLMV05]{FLMV05} L. Fortnow, R. Lipton, D. van Melkebeek and A. Viglas, \emph{Time-space lower bounds for satisfiability}, J. ACM \textbf{52} (2005), no.~6, 835--865. \href{https://doi.org/10.1145/1101821.1101822}{doi:10.1145/1101821.1101822}
\bibitem[For09]{For09} L. Fortnow, \emph{The status of the P versus NP problem}, Commun. ACM \textbf{52} (2009), no.~9, 78--86. \href{https://doi.org/10.1145/1562164.1562186}{doi:10.1145/1562164.1562186}
\bibitem[FSS84]{FSS84} M. Furst, J. B. Saxe and M. Sipser, \emph{Parity, circuits, and the polynomial-time hierarchy}, Math. Systems Theory \textbf{17} (1984), no.~1, 13--27. \href{https://doi.org/10.1007/BF01744431}{doi:10.1007/BF01744431}
\bibitem[Gas19]{Gas19} W. I. Gasarch, \emph{Guest column: the third P $=?$ NP poll}, ACM SIGACT News \textbf{50} (2019), no.~1, 38--59. \href{https://doi.org/10.1145/3319627.3319636}{doi:10.1145/3319627.3319636}
\bibitem[GHKK18]{GHKK18} A. Golovnev, E. A. Hirsch, A. Knop and A. S. Kulikov, \emph{On the limits of gate elimination}, J. Comput. System Sci. \textbf{96} (2018), 107--119. \href{https://doi.org/10.1016/j.jcss.2018.04.005}{doi:10.1016/j.jcss.2018.04.005}
\bibitem[GJ79]{GJ79} M. R. Garey and D. S. Johnson, \emph{Computers and Intractability: A Guide to the Theory of NP-Completeness}, W. H. Freeman, San Francisco, 1979.
\bibitem[GKST18]{GKST18} A. Golovnev, A. S. Kulikov, A. V. Smal and S. Tamaki, \emph{Gate elimination: circuit size lower bounds and \#SAT upper bounds}, Theoret. Comput. Sci. \textbf{719} (2018), 46--63. \href{https://doi.org/10.1016/j.tcs.2017.11.008}{doi:10.1016/j.tcs.2017.11.008}
\bibitem[Hak85]{Hak85} A. Haken, \emph{The intractability of resolution}, Theoret. Comput. Sci. \textbf{39} (1985), 297--308. \href{https://doi.org/10.1016/0304-3975(85)90144-6}{doi:10.1016/0304-3975(85)90144-6}
\bibitem[Has86]{Has86} J. H{\aa}stad, \emph{Almost optimal lower bounds for small depth circuits}, in: Proceedings of the 18th Annual ACM Symposium on Theory of Computing (STOC 1986), 1986, pp.~6--20. \href{https://doi.org/10.1145/12130.12132}{doi:10.1145/12130.12132}
\bibitem[Has98]{Has98} J. H{\aa}stad, \emph{The shrinkage exponent of de Morgan formulas is $2$}, SIAM J. Comput. \textbf{27} (1998), no.~1, 48--64. \href{https://doi.org/10.1137/S0097539794261556}{doi:10.1137/S0097539794261556}
\bibitem[Her14]{Her14} T. Hertli, \emph{3-SAT faster and simpler---Unique-SAT bounds for PPSZ hold in general}, SIAM J. Comput. \textbf{43} (2014), no.~2, 718--729. \href{https://doi.org/10.1137/120868177}{doi:10.1137/120868177}
\bibitem[Hir18]{Hir18} S. Hirahara, \emph{Non-Black-Box Worst-Case to Average-Case Reductions within NP}, in: 59th Annual IEEE Symposium on Foundations of Computer Science (FOCS 2018), 2018, pp.~247--258. \href{https://doi.org/10.1109/FOCS.2018.00032}{doi:10.1109/FOCS.2018.00032}
\bibitem[Hir22]{Hir22} S. Hirahara, \emph{NP-hardness of learning programs and partial MCSP}, in: 63rd Annual IEEE Symposium on Foundations of Computer Science (FOCS 2022), 2022, pp.~968--979. \href{https://doi.org/10.1109/FOCS54457.2022.00095}{doi:10.1109/FOCS54457.2022.00095}
\bibitem[HMPST93]{HMPST93} A. Hajnal, W. Maass, P. Pudl{\'a}k, M. Szegedy and G. Tur{\'a}n, \emph{Threshold circuits of bounded depth}, J. Comput. System Sci. \textbf{46} (1993), no.~2, 129--154. \href{https://doi.org/10.1016/0022-0000(93)90001-D}{doi:10.1016/0022-0000(93)90001-D}
\bibitem[HPV77]{HPV77} J. Hopcroft, W. Paul and L. Valiant, \emph{On Time Versus Space}, J. ACM \textbf{24} (1977), no.~2, 332--337. \href{https://doi.org/10.1145/322003.322015}{doi:10.1145/322003.322015}
\bibitem[HS65]{HS65} J. Hartmanis and R. E. Stearns, \emph{On the computational complexity of algorithms}, Trans. Amer. Math. Soc. \textbf{117} (1965), 285--306. \href{https://doi.org/10.1090/S0002-9947-1965-0170805-7}{doi:10.1090/S0002-9947-1965-0170805-7}
\bibitem[IKW02]{IKW02} R. Impagliazzo, V. Kabanets and A. Wigderson, \emph{In search of an easy witness: exponential time vs.\ probabilistic polynomial time}, J. Comput. System Sci. \textbf{65} (2002), no.~4, 672--694. \href{https://doi.org/10.1016/S0022-0000(02)00024-7}{doi:10.1016/S0022-0000(02)00024-7}
\bibitem[ILPS14]{ILPS14} R. Impagliazzo, S. Lovett, R. Paturi and S. Schneider, \emph{0-1 integer linear programming with a linear number of constraints}, preprint (2014), \href{https://arxiv.org/abs/1401.5512}{arXiv:1401.5512}.
\bibitem[ILW23]{ILW23} R. Ilango, J. Li and R. R. Williams, \emph{Indistinguishability obfuscation, range avoidance, and bounded arithmetic}, in: Proceedings of the 55th Annual ACM Symposium on Theory of Computing (STOC 2023), 2023, pp.~1076--1089. \href{https://doi.org/10.1145/3564246.3585187}{doi:10.1145/3564246.3585187}
\bibitem[IM02]{IM02} K. Iwama and H. Morizumi, \emph{An explicit lower bound of $5n-o(n)$ for Boolean circuits}, in: Mathematical Foundations of Computer Science 2002, Lecture Notes in Comput. Sci. \textbf{2420}, Springer, Berlin, 2002, pp.~353--364. \href{https://doi.org/10.1007/3-540-45687-2_29}{doi:10.1007/3-540-45687-2\_29}
\bibitem[Imp95]{Imp95} R. Impagliazzo, \emph{A personal view of average-case complexity}, in: Proceedings of the Tenth Annual Structure in Complexity Theory Conference, IEEE, 1995, pp.~134--147. \href{https://doi.org/10.1109/SCT.1995.514853}{doi:10.1109/SCT.1995.514853}
\bibitem[IP01]{IP01} R. Impagliazzo and R. Paturi, \emph{On the complexity of $k$-SAT}, J. Comput. System Sci. \textbf{62} (2001), no.~2, 367--375. \href{https://doi.org/10.1006/jcss.2000.1727}{doi:10.1006/jcss.2000.1727}
\bibitem[IPS13]{IPS13} R. Impagliazzo, R. Paturi and S. Schneider, \emph{A Satisfiability Algorithm for Sparse Depth Two Threshold Circuits}, in: 54th Annual IEEE Symposium on Foundations of Computer Science (FOCS 2013), 2013, pp.~479--488. \href{https://doi.org/10.1109/FOCS.2013.58}{doi:10.1109/FOCS.2013.58}
\bibitem[IPS97]{IPS97} R. Impagliazzo, R. Paturi and M. E. Saks, \emph{Size--Depth Tradeoffs for Threshold Circuits}, SIAM J. Comput. \textbf{26} (1997), no.~3, 693--707. \href{https://doi.org/10.1137/S0097539792282965}{doi:10.1137/S0097539792282965}
\bibitem[IPZ01]{IPZ01} R. Impagliazzo, R. Paturi and F. Zane, \emph{Which Problems Have Strongly Exponential Complexity?}, J. Comput. System Sci. \textbf{63} (2001), no.~4, 512--530. \href{https://doi.org/10.1006/jcss.2001.1774}{doi:10.1006/jcss.2001.1774}
\bibitem[IW97]{IW97} R. Impagliazzo and A. Wigderson, \emph{$P=BPP$ if $E$ requires exponential circuits: derandomizing the XOR lemma}, in: Proceedings of the 29th Annual ACM Symposium on Theory of Computing (STOC 1997), 1997, pp.~220--229. \href{https://doi.org/10.1145/258533.258590}{doi:10.1145/258533.258590}
\bibitem[JMV15]{JMV15} H. Jahanjou, E. Miles and E. Viola, \emph{Local Reductions}, in: Automata, Languages, and Programming (ICALP 2015), Part I, Lecture Notes in Comput. Sci. \textbf{9134}, Springer, Berlin, 2015, pp.~749--760. \href{https://doi.org/10.1007/978-3-662-47672-7_61}{doi:10.1007/978-3-662-47672-7\_61}
\bibitem[Juk10]{Juk10} S. Jukna, \emph{Representing (0, 1)-matrices by boolean circuits}, Discrete Math. \textbf{310} (2010), no.~1, 184--187. \href{https://doi.org/10.1016/j.disc.2009.07.011}{doi:10.1016/j.disc.2009.07.011}
\bibitem[Juk12]{Juk12} S. Jukna, \emph{Boolean Function Complexity}, Algorithms Combin. \textbf{27}, Springer, Heidelberg, 2012. \href{https://doi.org/10.1007/978-3-642-24508-4}{doi:10.1007/978-3-642-24508-4}
\bibitem[Kan82]{Kan82} R. Kannan, \emph{Circuit-size lower bounds and non-reducibility to sparse sets}, Inform. and Control \textbf{55} (1982), no.~1-3, 40--56. \href{https://doi.org/10.1016/S0019-9958(82)90382-5}{doi:10.1016/S0019-9958(82)90382-5}
\bibitem[Kar72]{Kar72} R. M. Karp, \emph{Reducibility among Combinatorial Problems}, in: Complexity of Computer Computations, R. E. Miller and J. W. Thatcher (eds.), Plenum Press, New York, 1972, pp.~85--103. \href{https://doi.org/10.1007/978-1-4684-2001-2_9}{doi:10.1007/978-1-4684-2001-2\_9}
\bibitem[KC00]{KC00} V. Kabanets and J.-Y. Cai, \emph{Circuit minimization problem}, in: Proceedings of the 32nd Annual ACM Symposium on Theory of Computing (STOC 2000), 2000, pp.~73--79. \href{https://doi.org/10.1145/335305.335314}{doi:10.1145/335305.335314}
\bibitem[KI04]{KI04} V. Kabanets and R. Impagliazzo, \emph{Derandomizing Polynomial Identity Tests Means Proving Circuit Lower Bounds}, Comput. Complexity \textbf{13} (2004), no.~1-2, 1--46. \href{https://doi.org/10.1007/s00037-004-0182-6}{doi:10.1007/s00037-004-0182-6}
\bibitem[Kor21]{Kor21} O. Korten, \emph{The hardest explicit construction}, in: 62nd Annual IEEE Symposium on Foundations of Computer Science (FOCS 2021), 2021, pp.~433--444. \href{https://doi.org/10.1109/FOCS52979.2021.00051}{doi:10.1109/FOCS52979.2021.00051}
\bibitem[Kra19]{Kra19} J. Kraj{\'\i}{\v c}ek, \emph{Proof Complexity}, Encyclopedia Math. Appl. \textbf{170}, Cambridge University Press, Cambridge, 2019. \href{https://doi.org/10.1017/9781108242066}{doi:10.1017/9781108242066}
\bibitem[KW16]{KW16} D. M. Kane and R. Williams, \emph{Super-linear gate and super-quadratic wire lower bounds for depth-two and depth-three threshold circuits}, in: Proceedings of the 48th Annual ACM Symposium on Theory of Computing (STOC 2016), 2016, pp.~633--643. \href{https://doi.org/10.1145/2897518.2897636}{doi:10.1145/2897518.2897636}
\bibitem[KW90]{KW90} M. Karchmer and A. Wigderson, \emph{Monotone Circuits for Connectivity Require Super-Logarithmic Depth}, SIAM J. Discrete Math. \textbf{3} (1990), no.~2, 255--265. \href{https://doi.org/10.1137/0403021}{doi:10.1137/0403021}
\bibitem[Lad75]{Lad75} R. E. Ladner, \emph{On the Structure of Polynomial Time Reducibility}, J. ACM \textbf{22} (1975), no.~1, 155--171. \href{https://doi.org/10.1145/321864.321877}{doi:10.1145/321864.321877}
\bibitem[Lev73]{Lev73} L. A. Levin, \emph{Universal sequential search problems}, Problemy Peredachi Informatsii \textbf{9} (1973), no.~3, 115--116 (in Russian); English translation in Problems Inform. Transmission \textbf{9} (1973), 265--266.
\bibitem[Li24]{Li24} Z. Li, \emph{Symmetric Exponential Time Requires Near-Maximum Circuit Size: Simplified, Truly Uniform}, in: Proceedings of the 56th Annual ACM Symposium on Theory of Computing (STOC 2024), 2024, pp.~2000--2007. \href{https://doi.org/10.1145/3618260.3649615}{doi:10.1145/3618260.3649615}
\bibitem[LST21]{LST21} N. Limaye, S. Srinivasan and S. Tavenas, \emph{Superpolynomial Lower Bounds Against Low-Depth Algebraic Circuits}, in: 62nd Annual IEEE Symposium on Foundations of Computer Science (FOCS 2021), 2021, pp.~804--814. \href{https://doi.org/10.1109/FOCS52979.2021.00083}{doi:10.1109/FOCS52979.2021.00083}
\bibitem[LY22]{LY22} J. Li and T. Yang, \emph{$3.1n-o(n)$ circuit lower bounds for explicit functions}, in: Proceedings of the 54th Annual ACM Symposium on Theory of Computing (STOC 2022), 2022, pp.~1180--1193. \href{https://doi.org/10.1145/3519935.3519976}{doi:10.1145/3519935.3519976}
\bibitem[MS01]{MS01} K. D. Mulmuley and M. Sohoni, \emph{Geometric complexity theory I: an approach to the P vs.\ NP and related problems}, SIAM J. Comput. \textbf{31} (2001), no.~2, 496--526. \href{https://doi.org/10.1137/S009753970038715X}{doi:10.1137/S009753970038715X}
\bibitem[MS08]{MS08} I. L. Markov and Y. Shi, \emph{Simulating Quantum Computation by Contracting Tensor Networks}, SIAM J. Comput. \textbf{38} (2008), no.~3, 963--981. \href{https://doi.org/10.1137/050644756}{doi:10.1137/050644756}
\bibitem[Mul12]{Mul12} K. D. Mulmuley, \emph{The GCT program toward the P vs.\ NP problem}, Commun. ACM \textbf{55} (2012), no.~6, 98--107. \href{https://doi.org/10.1145/2184319.2184341}{doi:10.1145/2184319.2184341}
\bibitem[MW20]{MW20} C. D. Murray and R. R. Williams, \emph{Circuit lower bounds for nondeterministic quasi-polytime from a new easy witness lemma}, SIAM J. Comput. \textbf{49} (2020), no.~5, STOC18-300--STOC18-322. \href{https://doi.org/10.1137/18M1195887}{doi:10.1137/18M1195887}
\bibitem[NW94]{NW94} N. Nisan and A. Wigderson, \emph{Hardness vs randomness}, J. Comput. System Sci. \textbf{49} (1994), no.~2, 149--167. \href{https://doi.org/10.1016/S0022-0000(05)80043-1}{doi:10.1016/S0022-0000(05)80043-1}
\bibitem[OAI26]{OAI26} OpenAI, \emph{The exponential time hypothesis is false}, preprint, 25 October 2026, \url{https://github.com/openai/math/blob/main/preprints/The-Exponential-Time-Hypothesis-is-False-October-25-2026/paper.pdf}.
\bibitem[OS18]{OS18} I. C. Oliveira and R. Santhanam, \emph{Hardness Magnification for Natural Problems}, in: 59th Annual IEEE Symposium on Foundations of Computer Science (FOCS 2018), 2018, pp.~65--76. \href{https://doi.org/10.1109/FOCS.2018.00016}{doi:10.1109/FOCS.2018.00016}
\bibitem[PPST83]{PPST83} W. J. Paul, N. Pippenger, E. Szemeredi and W. T. Trotter, \emph{On determinism versus non-determinism and related problems}, in: 24th Annual Symposium on Foundations of Computer Science (FOCS 1983), IEEE, 1983, pp.~429--438. \href{https://doi.org/10.1109/SFCS.1983.39}{doi:10.1109/SFCS.1983.39}
\bibitem[PPSZ05]{PPSZ05} R. Paturi, P. Pudl{\'a}k, M. E. Saks and F. Zane, \emph{An improved exponential-time algorithm for $k$-SAT}, J. ACM \textbf{52} (2005), no.~3, 337--364. \href{https://doi.org/10.1145/1066100.1066101}{doi:10.1145/1066100.1066101}
\bibitem[Pud94]{Pud94} P. Pudl{\'a}k, \emph{Communication in bounded depth circuits}, Combinatorica \textbf{14} (1994), no.~2, 203--216. \href{https://doi.org/10.1007/BF01215351}{doi:10.1007/BF01215351}
\bibitem[Rag88]{Rag88} P. Raghavan, \emph{Probabilistic construction of deterministic algorithms: Approximating packing integer programs}, J. Comput. System Sci. \textbf{37} (1988), no.~2, 130--143. \href{https://doi.org/10.1016/0022-0000(88)90003-7}{doi:10.1016/0022-0000(88)90003-7}
\bibitem[Raz87]{Raz87} A. A. Razborov, \emph{Lower bounds on the size of bounded depth circuits over a complete basis with logical addition}, Math. Notes Acad. Sci. USSR \textbf{41} (1987), no.~4, 333--338. \href{https://doi.org/10.1007/BF01137685}{doi:10.1007/BF01137685}
\bibitem[RM99]{RM99} R. Raz and P. McKenzie, \emph{Separation of the Monotone NC Hierarchy}, Combinatorica \textbf{19} (1999), no.~3, 403--435. \href{https://doi.org/10.1007/s004930050062}{doi:10.1007/s004930050062}
\bibitem[Ros08]{Ros08} B. Rossman, \emph{On the constant-depth complexity of $k$-clique}, in: Proceedings of the 40th Annual ACM Symposium on Theory of Computing (STOC 2008), 2008, pp.~721--730. \href{https://doi.org/10.1145/1374376.1374480}{doi:10.1145/1374376.1374480}
\bibitem[RR97]{RR97} A. A. Razborov and S. Rudich, \emph{Natural Proofs}, J. Comput. System Sci. \textbf{55} (1997), no.~1, 24--35. \href{https://doi.org/10.1006/jcss.1997.1494}{doi:10.1006/jcss.1997.1494}
\bibitem[RSW22]{RSW22} H. Ren, R. Santhanam and Z. Wang, \emph{On the range avoidance problem for circuits}, in: 63rd Annual IEEE Symposium on Foundations of Computer Science (FOCS 2022), 2022, pp.~640--650. \href{https://doi.org/10.1109/FOCS54457.2022.00067}{doi:10.1109/FOCS54457.2022.00067}
\bibitem[RTS00]{RTS00} J. Radhakrishnan and A. Ta-Shma, \emph{Bounds for Dispersers, Extractors, and Depth-Two Superconcentrators}, SIAM J. Discrete Math. \textbf{13} (2000), no.~1, 2--24. \href{https://doi.org/10.1137/S0895480197329508}{doi:10.1137/S0895480197329508}
\bibitem[San09]{San09} R. Santhanam, \emph{Circuit Lower Bounds for Merlin--Arthur Classes}, SIAM J. Comput. \textbf{39} (2009), no.~3, 1038--1061. \href{https://doi.org/10.1137/070702680}{doi:10.1137/070702680}
\bibitem[San10]{San10} R. Santhanam, \emph{Fighting Perebor: New and Improved Algorithms for Formula and QBF Satisfiability}, in: 51st Annual IEEE Symposium on Foundations of Computer Science (FOCS 2010), 2010, pp.~183--192. \href{https://doi.org/10.1109/FOCS.2010.25}{doi:10.1109/FOCS.2010.25}
\bibitem[Sav70]{Sav70} W. J. Savitch, \emph{Relationships between nondeterministic and deterministic tape complexities}, J. Comput. System Sci. \textbf{4} (1970), no.~2, 177--192. \href{https://doi.org/10.1016/S0022-0000(70)80006-X}{doi:10.1016/S0022-0000(70)80006-X}
\bibitem[Sch80]{Sch80} C. Schnorr, \emph{A $3n$-lower bound on the network complexity of Boolean functions}, Theoret. Comput. Sci. \textbf{10} (1980), no.~1, 83--92. \href{https://doi.org/10.1016/0304-3975(80)90074-2}{doi:10.1016/0304-3975(80)90074-2}
\bibitem[Sch99]{Sch99} U. Sch{\"o}ning, \emph{A probabilistic algorithm for $k$-SAT and constraint satisfaction problems}, in: 40th Annual Symposium on Foundations of Computer Science (FOCS 1999), IEEE, 1999, pp.~410--414. \href{https://doi.org/10.1109/SFFCS.1999.814612}{doi:10.1109/SFFCS.1999.814612}
\bibitem[SFM78]{SFM78} J. I. Seiferas, M. J. Fischer and A. R. Meyer, \emph{Separating Nondeterministic Time Complexity Classes}, J. ACM \textbf{25} (1978), no.~1, 146--167. \href{https://doi.org/10.1145/322047.322061}{doi:10.1145/322047.322061}
\bibitem[Sha49]{Sha49} C. E. Shannon, \emph{The Synthesis of Two-Terminal Switching Circuits}, Bell System Tech. J. \textbf{28} (1949), no.~1, 59--98. \href{https://doi.org/10.1002/j.1538-7305.1949.tb03624.x}{doi:10.1002/j.1538-7305.1949.tb03624.x}
\bibitem[Sip92]{Sip92} M. Sipser, \emph{The history and status of the P versus NP question}, in: Proceedings of the 24th Annual ACM Symposium on Theory of Computing (STOC 1992), 1992, pp.~603--618. \href{https://doi.org/10.1145/129712.129771}{doi:10.1145/129712.129771}
\bibitem[Smo87]{Smo87} R. Smolensky, \emph{Algebraic methods in the theory of lower bounds for Boolean circuit complexity}, in: Proceedings of the 19th Annual ACM Symposium on Theory of Computing (STOC 1987), 1987, pp.~77--82. \href{https://doi.org/10.1145/28395.28404}{doi:10.1145/28395.28404}
\bibitem[ST13]{ST13} K. Seto and S. Tamaki, \emph{A satisfiability algorithm and average-case hardness for formulas over the full binary basis}, Comput. Complexity \textbf{22} (2013), no.~2, 245--274. \href{https://doi.org/10.1007/s00037-013-0067-7}{doi:10.1007/s00037-013-0067-7}
\bibitem[Tal14]{Tal14} A. Tal, \emph{Shrinkage of de Morgan formulae by spectral techniques}, in: 55th Annual IEEE Symposium on Foundations of Computer Science (FOCS 2014), 2014, pp.~551--560. \href{https://doi.org/10.1109/FOCS.2014.65}{doi:10.1109/FOCS.2014.65}
\bibitem[Tar88]{Tar88} {\'E}. Tardos, \emph{The gap between monotone and non-monotone circuit complexity is exponential}, Combinatorica \textbf{8} (1988), no.~1, 141--142. \href{https://doi.org/10.1007/BF02122563}{doi:10.1007/BF02122563}
\bibitem[Tod91]{Tod91} S. Toda, \emph{PP is as hard as the polynomial-time hierarchy}, SIAM J. Comput. \textbf{20} (1991), no.~5, 865--877. \href{https://doi.org/10.1137/0220053}{doi:10.1137/0220053}
\bibitem[Val77]{Val77} L. G. Valiant, \emph{Graph-theoretic arguments in low-level complexity}, in: Mathematical Foundations of Computer Science 1977, Lecture Notes in Comput. Sci. \textbf{53}, Springer, Berlin, 1977, pp.~162--176. \href{https://doi.org/10.1007/3-540-08353-7_135}{doi:10.1007/3-540-08353-7\_135}
\bibitem[Val79a]{Val79a} L. G. Valiant, \emph{Completeness classes in algebra}, in: Proceedings of the 11th Annual ACM Symposium on Theory of Computing (STOC 1979), 1979, pp.~249--261. \href{https://doi.org/10.1145/800135.804419}{doi:10.1145/800135.804419}
\bibitem[Val79b]{Val79b} L. Valiant, \emph{The complexity of computing the permanent}, Theoret. Comput. Sci. \textbf{8} (1979), no.~2, 189--201. \href{https://doi.org/10.1016/0304-3975(79)90044-6}{doi:10.1016/0304-3975(79)90044-6}
\bibitem[vzGG13]{vzGG13} J. von zur Gathen and J. Gerhard, \emph{Modern Computer Algebra}, third edition, Cambridge University Press, Cambridge, 2013. \href{https://doi.org/10.1017/CBO9781139856065}{doi:10.1017/CBO9781139856065}
\bibitem[Wig19]{Wig19} A. Wigderson, \emph{Mathematics and Computation: A Theory Revolutionizing Technology and Science}, Princeton University Press, Princeton, 2019.
\bibitem[Wil05]{Wil05} R. Williams, \emph{A new algorithm for optimal 2-constraint satisfaction and its implications}, Theoret. Comput. Sci. \textbf{348} (2005), no.~2-3, 357--365. \href{https://doi.org/10.1016/j.tcs.2005.09.023}{doi:10.1016/j.tcs.2005.09.023}
\bibitem[Wil13]{Wil13} R. Williams, \emph{Improving exhaustive search implies superpolynomial lower bounds}, SIAM J. Comput. \textbf{42} (2013), no.~3, 1218--1244. \href{https://doi.org/10.1137/10080703X}{doi:10.1137/10080703X}
\bibitem[Wil14]{Wil14} R. Williams, \emph{Nonuniform ACC circuit lower bounds}, J. ACM \textbf{61} (2014), no.~1, 1--32. \href{https://doi.org/10.1145/2559903}{doi:10.1145/2559903}
\bibitem[Wil14b]{Wil14b} R. Williams, \emph{New algorithms and lower bounds for circuits with linear threshold gates}, in: Proceedings of the 46th Annual ACM Symposium on Theory of Computing (STOC 2014), 2014, pp.~194--202. \href{https://doi.org/10.1145/2591796.2591858}{doi:10.1145/2591796.2591858}
\bibitem[Wil14c]{Wil14c} R. Williams, \emph{Algorithms for Circuits and Circuits for Algorithms}, in: 29th IEEE Conference on Computational Complexity (CCC 2014), 2014, pp.~248--261. \href{https://doi.org/10.1109/CCC.2014.33}{doi:10.1109/CCC.2014.33}
\bibitem[Wil16]{Wil16} R. R. Williams, \emph{Natural Proofs versus Derandomization}, SIAM J. Comput. \textbf{45} (2016), no.~2, 497--529. \href{https://doi.org/10.1137/130938219}{doi:10.1137/130938219}
\bibitem[Wil25]{Wil25} R. R. Williams, \emph{Simulating Time with Square-Root Space}, in: Proceedings of the 57th Annual ACM Symposium on Theory of Computing (STOC 2025), 2025, pp.~13--23. \href{https://doi.org/10.1145/3717823.3718225}{doi:10.1145/3717823.3718225}
\bibitem[Woe16]{Woe16} G. J. Woeginger, \emph{The P-versus-NP page}, \url{https://wscor.win.tue.nl/woeginger/P-versus-NP.htm}, last updated 2016.
\bibitem[Yan91]{Yan91} M. Yannakakis, \emph{Expressing combinatorial optimization problems by linear programs}, J. Comput. System Sci. \textbf{43} (1991), no.~3, 441--466. \href{https://doi.org/10.1016/0022-0000(91)90024-Y}{doi:10.1016/0022-0000(91)90024-Y}
\bibitem[Yao85]{Yao85} A. C.-C. Yao, \emph{Separating the polynomial-time hierarchy by oracles}, in: 26th Annual Symposium on Foundations of Computer Science (FOCS 1985), IEEE, 1985, pp.~1--10. \href{https://doi.org/10.1109/SFCS.1985.49}{doi:10.1109/SFCS.1985.49}
\bibitem[Zak83]{Zak83} S. {\v Z}{\'a}k, \emph{A Turing machine time hierarchy}, Theoret. Comput. Sci. \textbf{26} (1983), no.~3, 327--333. \href{https://doi.org/10.1016/0304-3975(83)90015-4}{doi:10.1016/0304-3975(83)90015-4}
\bibitem[Zoo]{Zoo} S. Aaronson et al., \emph{The Complexity Zoo}, \url{https://complexityzoo.net/Complexity_Zoo}.
\end{thebibliography}
\endgroup
